% =============================================================================
% FST-Nash: Spieltheoretische Diagnostik von Chaperon-Systemen
% Version 1.4 -- Zenodo Preprint
% Autor: Lukas Geiger, unabhängiger Forscher, Bernau
% =============================================================================
\documentclass[12pt,a4paper]{article}

\usepackage[utf8]{inputenc}
\usepackage[T1]{fontenc}
\usepackage[ngerman]{babel}
\usepackage{lmodern}
\usepackage{microtype}
\usepackage{amsmath,amssymb,amsthm}
\usepackage{graphicx}
\usepackage{booktabs}
\usepackage{multirow}
\usepackage{array}
\usepackage{longtable}
\usepackage{tabularx}
\usepackage{needspace}
\usepackage[margin=2.5cm]{geometry}
\usepackage{xcolor}
\usepackage{xurl}
\input{glyphtounicode}
\pdfgentounicode=1
\usepackage[unicode=true,colorlinks=true,linkcolor=blue,citecolor=blue,urlcolor=blue]{hyperref}
\hypersetup{
  pdftitle={Konstruktionsbedingte Potentialspiel-Diagnostik für Chaperon-Modelle},
  pdfauthor={Lukas Geiger},
  pdfsubject={Functional Stability Theory; Chaperon-Systeme; Potentialspiel-Diagnostik},
  pdfkeywords={Chaperon-Systeme, Potentialspiel, Nash-Gleichgewicht, Symmetriebrechung, Proteinfaltung, Nichtgleichgewichtsthermodynamik, Spieltheorie, Faltungswechsel, Functional Stability Theory},
  pdflang={de-DE},
  bookmarksnumbered=true,
  pdfpagemode=UseOutlines
}

\newtheorem{proposition}{Proposition}
\newtheorem{definition}{Definition}
\renewcommand{\proofname}{Beweis}
\newcolumntype{L}[1]{>{\raggedright\arraybackslash}p{#1}}
\newcolumntype{Y}{>{\raggedright\arraybackslash}X}
\emergencystretch=3em

\title{Konstruktionsbedingte Potentialspiel-Diagnostik\\
für Chaperon-Modelle}

\author{Lukas Geiger\\
\small Unabhängiger Forscher, Bernau\\
\small Code: \href{https://github.com/research-line/fst-nash}{research-line/fst-nash}}

\date{Oktober 2026}

\begin{document}
\maketitle

% =============================================================================
\begin{abstract}
Wir untersuchen $2\times 2$-Normalformspiele, die aus ausgewählten
Chaperon--Substrat- und Faltungswechsel-Beispielen konstruiert wurden,
mit dem exakten Potentialspiel-Test (PG) von Monderer \& Shapley (1996).
Unter der hier verwendeten Modellierungskonvention wird ein geteilter
Interaktionskontrast als PG und ein asymmetrischer Kontrast als non-PG
kodiert. Diese Konvention ist durch Xus Symmetrierahmen (2022)
motiviert, wird aber nicht aus einem Chaperon-Reaktionsnetz hergeleitet
und belegt daher keine $\mathrm{S4}\iff\mathrm{PG}$-Äquivalenz.

Die Tabelle mit 16 Fällen ist folglich ein \emph{Konstruktionsledger},
kein unabhängig validierter biologischer Atlas. XCL1/\allowbreak
Lymphotactin ermöglicht eine interne Rekonstruktion der Salzverschiebung
($23{,}7\,\%$~Ltn40 bei $150$\,mM~NaCl), da derselbe Datensatz von Tyler et~al.\ (2012)
sowohl die Kalibrierung als auch den Vergleich bereitstellt. Das Thermosom-Beispiel
zeigt, dass dasselbe biologische System unter verschiedenen
Spielkonstruktionen entgegengesetzte PG-Labels erhalten kann. Ein
Validierungsledger mit fünf Achsen ergibt $0/8$ vollständige Pässe.
Nash-Gleichgewichtsidentitäten und Sensitivitätsgrenzen sind daher
Kandidaten für künftige vorregistrierte Tests, keine validierten
Vorhersagen.

\medskip\noindent
\textbf{Bezug zur Vorarbeit.}
Die vorliegende Arbeit ersetzt Abschnitt~3 (\glqq Game-Theoretic Stability\grqq{}) des
Übersichtspapiers \emph{Functional Stability Theory~III: Biological
Stability and Nash Frustration} (Geiger 2026, DOI:
10.5281/zenodo.20130573). Die dort berichtete Nash-Frustration-Korrelation
($\rho_S = 0.44$, $p = 0.033$, $n = 24$) ist unter der früheren
Minimierungs- und Schrittweitenkonstruktion nicht trennscharf: Für
eine positiv definite Hesse-Matrix folgt $\rho(I-\eta H)<1$, sofern
$0<\eta<2/\lambda_{\max}(H)$ gilt. Diese Resultate werden daher hier
nicht als biologische Evidenz weitergeführt; der vorliegende Rahmen
verwendet eine andere Konstruktion.
\end{abstract}

\clearpage
\begingroup
\small
\setlength{\parskip}{0pt}
\tableofcontents
\endgroup
\clearpage


% =============================================================================
\section{Einleitung}
\label{sec:intro}

Molekulare Chaperone unterstützen Proteinfaltung durch sehr verschiedene
Mechanismen, von passiver Holdase-Aktivität bis zu ATP-getriebenen
Konformationszyklen. Xu~\cite{Xu2022}, aufbauend auf Goloubinoff et~al.~\cite{Goloubinoff2018},
klassifizierte Chaperon-Mechanismen über
vier Symmetriebedingungen an der Ratenmatrix von Chaperon--Substrat-
Interaktionen:
\begin{itemize}
  \item[\textbf{S1}] detaillierte Bilanz von Bindung und Entbindung
  \item[\textbf{S2}] substratunabhängige Dissoziation
  \item[\textbf{S3}] substratunabhängige ATPase-Rate
  \item[\textbf{S4}] Reversibilität: Jeder geschlossene Chaperon-Zustand
    hat einen reversiblen Pfad zu einem offenen Zustand, ohne chemische
    Energie zu verbrauchen
\end{itemize}
Wenn alle vier Bedingungen gelten, befindet sich das System im
Gleichgewicht. Der Bruch einer beliebigen Bedingung führt zu einem
Nichtgleichgewichtssystem. Wir kürzen diese Klassen im Folgenden als
Goldilocks~(GG) und Nicht-Goldilocks~(NGG) ab; diese Terminologie ist
unsere, nicht die von Xu. Den Ausdruck \glqq Goloubinoff-Landschaft\grqq{}
verwenden wir als Kurzform für das Forschungsprogramm zu Gleichgewicht
und Nichtgleichgewicht bei Chaperonen, das von Goloubinoff et~al.~\cite{Goloubinoff2018}
begründet und von Xu~\cite{Xu2022} in die
S1--S4-Bedingungen formalisiert wurde. Xu unterscheidet Hsp70 und Hsp90
bereits anhand der verschiedenen Teilmengen gebrochener
Symmetriebedingungen und ihrer unterschiedlichen
Nichtgleichgewichtsmechanismen. Unsere konkrete Fragestellung lautet, ob explizit
festgelegte Spielrepräsentationen solche Symmetriemuster nachvollziehbar
rekodieren können.

\subsection{Beitrag}

Wir wenden die \emph{Potentialspiel}-Eigenschaft (PG) von Monderer \&
Shapley~\cite{MondererShapley1996} auf explizit konstruierte
$2\times 2$-Normalformspiele an. Der Vierzyklus-Test prüft eine
festgelegte Auszahlungsrepräsentation; er erlaubt jedoch keine unabhängige
Bestimmung der S4-Bedingung von Xu. Wir verwenden die folgenden vorgeschlagenen Labels
als beschreibende Unterteilung:
\begin{enumerate}
  \item \textbf{Kinetisches NGG} (S3 gebrochen, S4 als intakt behandelt;
    mit geteiltem Kontrast kodiert, daher PG per Konstruktion):
    vorgeschlagen für Routen, deren Raten substratabhängig sind, für die
    aber kein gerichteter Zyklus angenommen wird, der das Gleichgewicht
    verschiebt.
  \item \textbf{Thermodynamisches NGG} (S3 und S4 als gebrochen behandelt;
    mit asymmetrischen Kontrasten kodiert, daher non-PG per Konstruktion):
    vorgeschlagen für Routen, in denen ein gerichteter thermodynamischer
    Zyklus angenommen wird, der das Gleichgewicht verschiebt; die
    vorliegenden Modelle berechnen keine Zyklusaffinitäten.
\end{enumerate}

Wir dokumentieren 16~Modellkonstruktionen, darunter vier nachträgliche
Konventions-Konsistenzfälle, rekonstruieren die kalibrierte
XCL1-Salzantwort und zeigen die Konstruktionsbedingtheit anhand der
Thermosom-Meta-Demonstration.

\subsection{Bezug zur Vorarbeit}

Die vorliegende Arbeit ersetzt die Nash-Frustration-Analyse in FST-III~\cite{GeigerFSTIII2025}. Dort wurde jeder Rest als Spieler modelliert,
der Rückgrat-Diederwinkel $(\varphi_i, \psi_i)$ auf dem Torus~$T^2$
wählt, mit paarweisen Kosinus-Kopplungen. Der Spektralradius
$\rho(J) < 1$ der Jacobi-Matrix $J = I - \eta H$ wurde als
\glqq Nash-Stabilität\grqq{} interpretiert. Eine Zufallsparameter-Kontrolle
(Abschnitt~\ref{sec:tautology}) zeigte jedoch, dass für eine positiv
definite Hesse-Matrix $\rho(I-\eta H)<1$ unter der
Schrittweitenbedingung $0<\eta<2/\lambda_{\max}(H)$ folgt. Der Test ist
unter dieser Minimierungs- und Schrittweitenkonstruktion daher nicht
trennscharf. Frühere Resultate einschließlich der
Spearman-Korrelation $\rho_S = 0.44$ gegen B-Faktoren werden hier nicht
als biologische Evidenz weitergeführt.

Der vorliegende Rahmen gibt den kontinuierlichen Diederwinkel-Ansatz
vollständig auf und formuliert Chaperon--Substrat-Interaktionen
stattdessen als diskrete $2\times 2$-Normalformspiele, in denen der PG-Test eine
wohldefinierte Prüfung der festgelegten Auszahlungsrepräsentation ist.


% =============================================================================
\section{Theorie: Vierzyklus-Test und S4-motivierte Konstruktion}
\label{sec:theory}

\begin{definition}[Potentialspiel, Monderer \& Shapley 1996]
Ein endliches $n$-Personen-Spiel ist ein \emph{Potentialspiel}, wenn es
eine Funktion $\Phi: S_1 \times \cdots \times S_n \to \mathbb{R}$ gibt,
sodass für jeden Spieler~$k$, jedes Strategieprofil $s_{-k}$ und jedes
Paar von Strategien $s_k, s_k'$ gilt:
\begin{equation}
  u_k(s_k', s_{-k}) - u_k(s_k, s_{-k})
  = \Phi(s_k', s_{-k}) - \Phi(s_k, s_{-k}).
  \label{eq:pg_def}
\end{equation}
\end{definition}

Für ein $2\times 2$-Spiel mit Spielern $H$ (Chaperon) und $S$
(Substrat), Strategien $\{0, 1\}$ und Auszahlungsmatrizen $u_H$, $u_S$
reduziert sich die PG-Bedingung auf einen einzigen skalaren Test:

\begin{proposition}[Vierzyklus-Test für $2\times 2$-Spiele]
\label{prop:four_cycle}
Definiere den \emph{Interaktionsindex} für jeden Spieler~$k$:
\begin{equation}
  I_k = u_k[0,0] - u_k[0,1] - u_k[1,0] + u_k[1,1].
  \label{eq:interaction}
\end{equation}
Das Spiel ist genau dann ein Potentialspiel, wenn $I_H = I_S$.
\end{proposition}

\begin{proof}
Betrachte die beiden unilateralen Pfade von $(0,0)$ nach $(1,1)$. Auf
dem ersten Pfad wechselt $H$ die Zeile und anschließend $S$ die Spalte;
auf dem zweiten Pfad wechselt zuerst $S$ die Spalte und dann $H$ die
Zeile. Ein skalares exaktes Potential ist genau dann pfadunabhängig,
wenn
\begin{align}
 &[u_H(1,0)-u_H(0,0)]+[u_S(1,1)-u_S(1,0)] \nonumber\\
 ={}&[u_S(0,1)-u_S(0,0)]+[u_H(1,1)-u_H(0,1)].
\end{align}
Durch Umstellen folgt $I_H=I_S$, also die Vierzyklusbedingung für diese
Orientierung. Die Umkehrung folgt durch Konstruktion von $\Phi$ entlang
eines der beiden Pfade~\cite{MondererShapley1996}.
\end{proof}

\subsection{Durch Xus S4 motivierte Konstruktionsregel}

Xus vierte Symmetriebedingung~(S4)~\cite{Xu2022} fordert, dass jeder
geschlossene Chaperon-Zustand einen reversiblen Pfad zu einem offenen
Zustand hat, ohne chemische Energie zu verbrauchen -- also dass das
Chaperon--Substrat-Reaktionsnetzwerk mikroskopisch reversibel ist.

In den vorliegenden Modellen wird ein als S4-intakt behandelter Fall mit
einem geteilten Interaktionskontrast dargestellt; dadurch gilt
$I_H=I_S$ konstruktionsbedingt. Ein als S4-gebrochen behandelter Fall
wird mit asymmetrischen Kontrasten dargestellt; dadurch gilt im
Allgemeinen $I_H\neq I_S$. Dies sind Modellierungsregeln, keine
Herleitung aus Übergangsraten und kein Beweis biologischer Äquivalenz.

\textbf{Konstruktionsabhängigkeit.} Potentialität ist eine Eigenschaft
des festgelegten Spiels. Spielerpartition, Grobkörnung,
Nutzendefinition (\emph{utility}) und kardinale Skala können das PG-Ergebnis ändern.
Ein grobgekörntes Vierzustandsspiel kann zudem dissipative Zyklen
verbergen. PG impliziert daher weder biologische Reversibilität noch S4,
solange keine vollständige Netzwerk-zu-Spiel-Abbildung und unabhängig
bestimmte Zyklusaffinitäten vorliegen. Die Thermosom-Meta-Demonstration
(Abschnitt~\ref{sec:thermosome}) zeigt diese Abhängigkeit.


% =============================================================================
\section{Methoden}
\label{sec:methods}

\subsection{Spielkonstruktion}
\label{sec:game_construction}

Jedes Chaperon--Substrat-System wird als $2\times 2$-Normalformspiel
modelliert:
\begin{itemize}
  \item \textbf{Spieler~$H$} (Chaperon): Strategien
    $\{\text{Inaktiv/Ruhezustand}, \text{Aktiv/Gebunden}\}$
  \item \textbf{Spieler~$S$} (Substrat): Strategien
    $\{\text{Entfaltet/Frei}, \text{Gefaltet/Gebunden}\}$
\end{itemize}

Die Auszahlungsmatrix für Spieler~$k \in \{H, S\}$ ist:
\begin{equation}
  u_k = \begin{pmatrix}
    u_k[0,0] & u_k[0,1] \\
    u_k[1,0] & u_k[1,1]
  \end{pmatrix},
\end{equation}
wobei größere Einträge rechnerisch als bevorzugte Modellscores behandelt
werden. Einige Terme sind energieabgeleitet und in kcal/mol angegeben;
andere sind Transformationen kinetischer Verhältnisse oder Annahmen.
Die gegenwärtigen Fallmodelle liefern keine universelle
thermodynamische Abbildung von $K_D$, $k_\text{cat}$ und
Konformationsenergien auf eine gemeinsame kardinale Utility-Skala.
Exact-PG-Differenzen in kcal/mol sind daher konstruktionsbezogene
Scorekontraste, keine direkten physikalischen Messungen. Ein
thermodynamisches Modell mit eigenständigem wissenschaftlichem Anspruch benötigt stattdessen gepaarte
Vorwärts-/Rückwärtsraten oder ausdrücklich definierte Utilities wie
$u=-G$ bei fester Temperatur und festem Standardzustand.

Die Kopplungsenergie $\Delta_k$ beschreibt die Änderung der Auszahlung
von Spieler~$k$, wenn der andere Spieler die Strategie wechselt:
\begin{equation}
  \Delta_k = u_k[\text{gekoppelt}] - u_k[\text{ungekoppelt}].
\end{equation}

\subsection{Vierzyklus-Test (PG-Test)}

Für jedes kalibrierte System berechnen wir $I_H$ und $I_S$ gemäß
Gl.~\eqref{eq:interaction} und werten aus:
\begin{equation}
  \text{PG-Verletzung} = |I_H - I_S|.
\end{equation}
Liegt die Verletzung unterhalb einer numerischen Schwelle
($\varepsilon = 10^{-10}$), wird das festgelegte Spiel als exaktes PG
klassifiziert. Diese numerische Identität ist keine empirische
S4-Entscheidung und pflanzt Messunsicherheit nicht fort.

\subsection{Suche nach Nash-Gleichgewichten}

Reine Nash-Gleichgewichte werden identifiziert, indem für alle vier
Strategieprofile einseitige Abweichungsanreize geprüft werden. Ein Profil
$(s_H^*, s_S^*)$ ist ein reines NE, wenn gilt:
\begin{align}
  u_H(s_H^*, s_S^*) &\geq u_H(s_H', s_S^*) \quad \forall\, s_H', \\
  u_S(s_H^*, s_S^*) &\geq u_S(s_H^*, s_S') \quad \forall\, s_S'.
\end{align}

\subsection{Regime-Klassifikation}

Auf Basis des PG-Tests und der Symmetriebedingungen von Xu~\cite{Xu2022}:
\begin{enumerate}
  \item \textbf{GG} (Gleichgewicht): S1--S4 als intakt behandelt; die
    gewählte Repräsentation mit geteiltem Kontrast ist PG; kein ATP oder
    keine modellierte Substratstimulation
  \item \textbf{Kinetisches NGG}: S3 gebrochen, S4 als intakt behandelt;
    die gewählte Repräsentation mit geteiltem Kontrast ist PG
  \item \textbf{Thermodynamisches NGG}: S3 und S4 als gebrochen
    behandelt; die gewählte Repräsentation mit asymmetrischem Kontrast
    ist non-PG
  \item \textbf{Spezialfälle}: nichtstandardisierte $2\times 2$-
    Struktur (3-Spieler, oszillatorisch, reguliert)
\end{enumerate}

\subsection{Symmetrie-Mapping}

Für jedes System bewerten wir S1--S4 anhand publizierter experimenteller
Daten:
\begin{itemize}
  \item S1: Erfüllt das System detaillierte Bilanz für Bindung und
    Entbindung?
  \item S2: Ist die Dissoziation substratunabhängig?
  \item S3: Ist die ATPase-Rate substratunabhängig?
  \item S4: Erfüllt die zitierte biologische Route Xus
    Reversibilitätsbedingung? Dieses Label wird getrennt vom
    Spielergebnis dokumentiert; die vorliegende Konstruktion kodiert ein
    S4-intaktes Label mit geteiltem Kontrast und damit mit $I_H=I_S$.
\end{itemize}


% =============================================================================
\section{Ergebnisse}

% --- Result 1 ----------------------------------------------------------------
\subsection{Konstruktionsbedingter Regime-Vorschlag}
\label{sec:regime_trinity}

Die 16~Spielkonstruktionen lassen sich in drei vorgeschlagene Labels
ordnen (Tabelle~\ref{tab:atlas}). Dies ist eine beschreibende Gruppierung
der gewählten Repräsentationen, keine unabhängig gelernte biologische
Unterteilung:

\begin{enumerate}
  \item \textbf{Gleichgewicht (GG):} 6~Systeme. PG $=$ True, kein
    S3-Bruch. Dazu gehören XCL1 (metamorpher Faltungswechsel),
    Prefoldin (Holdase) und 4~nachträgliche
    Konventions-Konsistenzsysteme (DnaJ, Hsp26, Hsp33, Hsc70).

  \item \textbf{Kinetisches NGG:} 2~Systeme (Hsp70/DnaJ, SecA).
    Xu beschreibt Hsp70-vermittelte Faltung als S3-Bruch. Seine
    notwendige Bedingung für einen S4-Bruch verlangt mindestens zwei
    sequenzielle ATP-Hydrolysen pro Zyklus; diese fehlen in der hier
    betrachteten Hsp70-Route. Das biologische S4-Label ist daher mit Xu
    vereinbar. Davon getrennt folgt PG$=$True in unserer Matrix aus dem
    angenommenen geteilten Interaktionskontrast; es ist kein
    unabhängiger Test dieses Labels.

  \item \textbf{Thermodynamisches NGG:} 4~Systeme (Hsp90/Cdc37,
    GroEL/GroES, ClpB/Hsp104, p97/VCP). PG $=$ False mit konstruierten
    Scoreunterschieden von $0{,}5$--$2{,}5$. Diese Systeme werden als Bruch von S3 und S4 dargestellt. Die
    motivierende Hypothese lautet, dass Co-Chaperone, Mehrringarchitekturen
    oder Domänenasymmetrien gerichtete thermodynamische Zyklen erzeugen, die
    das Gleichgewicht verschieben; die vorliegenden $2\times 2$-Scores
    kodieren diese Annahme und berechnen oder prüfen keine Zyklusaffinitäten. Bei GroEL erfordert die
    Substratfreisetzung sequenzielle ATP-Hydrolyse im trans-Ring und
    verletzt damit Xus S4~\cite{Xu2022}.
\end{enumerate}

Vier zusätzliche Systeme (Trigger Factor, KaiB/KaiC, Hsp70$\to$Hsp90-Handoff,
MAD2) liegen aufgrund nichtstandardisierter Spielstrukturen
außerhalb der drei Standardregime (ribosomengekoppelt, oszillatorisch,
3-Spieler oder reguliert metamorph).


% --- Result 2 ----------------------------------------------------------------
\clearpage
\subsection{Konstruktionsledger}
\label{sec:atlas}

\begingroup
\footnotesize
\setlength{\tabcolsep}{2.8pt}
\renewcommand{\arraystretch}{1.12}
\begin{longtable}{@{}rL{0.23\textwidth}cccL{0.13\textwidth}L{0.33\textwidth}@{}}
\caption{Konstruktionsledger: 16 festgelegte Spielrepräsentationen,
geprüft mit dem Exact-PG-Test. Die letzte Spalte dokumentiert
Eingabestatus und die Konstruktion, die PG oder non-PG bestimmt. Keine
Zeile validiert unabhängig ein biologisches Label
(siehe Abschnitt~\ref{sec:thermosome}). NE $=$ Anzahl reiner
Nash-Gleichgewichte im festgelegten Spiel.}
\label{tab:atlas}
\\
\toprule
\# & System & ATP? & PG & NE & Regime & Evidenz / PG-Grundlage \\
\midrule
\endfirsthead
\toprule
\# & System & ATP? & PG & NE & Regime & Evidenz / PG-Grundlage \\
\midrule
\endhead
\bottomrule
\endfoot
\bottomrule
\endlastfoot
\multicolumn{7}{@{}l}{\textit{Regime 1 -- Gleichgewicht (GG)}} \\
1  & XCL1 (metamorph) & Nein & $\checkmark$ & \textbf{2} & GG
   & Symmetrisches Spiel per Konstruktion; kalibriert. \\
2  & Prefoldin (Holdase) & Nein & $\checkmark$* & \textbf{2}* & GG*
   & Geschätzte Parameter ($f{=}0{,}5$); nicht unabhängig kalibriert. \\
3  & DnaJ allein & Nein & $\checkmark$ & 1 & GG
   & Geteilte-$\Delta$-Konvention; nachträglicher Konsistenzfall. \\
4  & Hsp26 sHsp & Nein & $\checkmark$ & 1 & GG
   & Geteilte-$\Delta$-Konvention; nachträglicher Konsistenzfall. \\
5  & Hsp33 redox & Nein & $\checkmark$ & 1 & GG
   & Geteilte-$\Delta$-Konvention; nachträglicher Konsistenzfall. \\
6  & Hsc70 allein & Basal & $\checkmark$ & 1 & GG
   & Geteilte-$\Delta$-Konvention; nachträglicher Konsistenzfall. \\
\midrule
\multicolumn{7}{@{}l}{\textit{Regime 2 -- Kinetisches NGG}} \\
7  & Hsp70/DnaJ & Ja & $\checkmark$ & 1 & kin.\ NGG
   & Geteiltes $\Delta$ aus experimentellen $K_D$-Werten; kalibriert. \\
8  & SecA & Ja & $\checkmark$ & 1 & kin.\ NGG
   & Geteiltes $\Delta$ aus experimentellen $K_D$-Werten; kalibriert. \\
\midrule
\multicolumn{7}{@{}l}{\textit{Regime 3 -- Thermodynamisches NGG}} \\
9  & Hsp90/Cdc37 & Ja & $\times$ & 1 & therm.\ NGG
   & Asymmetrisches $\Delta$ (kalibriert/geschätzt); kalibriert. \\
10 & GroEL/GroES & Ja & $\times$ & 1 & therm.\ NGG
   & Experimentell motivierte Ringasymmetrie; kalibriert. \\
11 & ClpB/Hsp104 & Ja & $\times$ & 1 & therm.\ NGG
   & Threading-Asymmetrie; kalibriert. \\
12 & p97/VCP & Ja & $\times$ & 1 & therm.\ NGG
   & Geschätzte D1/D2-Asymmetrie; kalibriert. \\
\midrule
\multicolumn{7}{@{}l}{\textit{Spezialfälle (nichtstandardisiertes Regime)}} \\
13 & Trigger Factor & Nein & $\times$ & 1 & non-GG
   & Ribosomenasymmetrie mit experimentellem $K_D$; kalibriert. \\
14 & KaiB/KaiC & Indirekt & $\times$ & 1 & non-GG
   & Oszillatorisches System mit experimenteller Kinetik; kalibriert. \\
15 & Hsp70$\to$Hsp90 (3-Spieler) & Ja & $\times$ & 1 & non-PG
   & Drei-Spieler-Modell, kein $2\times 2$-Spiel; kalibriert. \\
16 & MAD2 (reguliert metamorph) & Nein$^\ddagger$ & $\times$ & 1 & non-PG
   & TRIP13-regulierte Asymmetrie; kalibriert. \\
\end{longtable}

\noindent{\footnotesize
\textbf{Ehrlichkeitshinweis.} Alle PG$=$True-Systeme haben PG als
algebraische Folge der Modellierungskonvention: XCL1 aus einem
symmetrischen Spiel (identische Spieler), Hsp70/DnaJ und SecA aus der
Annahme $\Delta_H = \Delta_S$ (experimentelle $K_D$-Werte, aber
symmetrische Kopplung wird \emph{angenommen}) und nachträgliche Fälle
aus der Geteilte-$\Delta$-Konvention. Alle non-PG-Systeme haben non-PG aus
$\Delta_H \neq \Delta_S$ oder aus asymmetrischer Spielstruktur.
PG/non-PG ist damit eine \emph{Konventionsentscheidung während der
Modellierung}, kein emergentes empirisches Resultat. * Prefoldin basiert
auf geschätzten Parametern ($K_D$, $\Delta G_\text{conf}$ unveröffentlicht).
$^\ddagger$ MAD2 hat kein direktes ATP; TRIP13 AAA+ ATPase setzt MAD2
zurück.}
\endgroup


% --- Result 3 ----------------------------------------------------------------
\subsection{XCL1-Salzverschiebung: interne kalibrierte Rekonstruktion}
\label{sec:xcl1}

XCL1/Lymphotactin ist ein metamorphes Protein, das zwischen einer
Chemokin-Faltung (Ltn10) und einer $\beta$-Faltblatt-Faltung (Ltn40)
interkonvertiert, mit $K_\text{eq} = 1{,}2$ bei $25\,{}^\circ\text{C}$ und
$0$\,mM~NaCl als Kalibrierungsanker des Projekts. Die internen
Aufzeichnungen verweisen sowohl auf Tuinstra et~al.~\cite{Tuinstra2008}
als auch auf Tyler et~al.~\cite{Tyler2012}; die genaue Zuordnung dieses
Werts zum Primärvolltext bleibt ungeprüft. Wir modellieren XCL1 als symmetrisches
$2\times 2$-Koordinationsspiel zwischen zwei identischen Monomeren, mit
den Strategien $\{\text{Chemokin}, \text{Beta-Faltblatt}\}$.

Das Spiel ergibt 2~reine NE: (Chem, Chem) $=$ Ltn10 und
(Beta, Beta) $=$ Ltn40. Damit wird die metamorphe Bistabilität als
Koordinationsgleichgewicht dargestellt.

\subsubsection{Salzabhängige Populationsverschiebung}

Tyler et~al.~\cite{Tyler2012} zeigten, dass NaCl die Population von
${\sim}55\,\%$~Ltn40 ($0$\,mM) auf ${\sim}24\,\%$~Ltn40 ($150$\,mM)
verschiebt, vermittelt durch Chloridbindung an Arg23/Arg43
($K_d \approx 16$\,mM aus CSP-Titration). Wir modellieren dies als
Chloridbindungs-Sättigung:

\begin{equation}
  f(\text{Cl}^-) = \frac{[\text{Cl}^-]}{[\text{Cl}^-] + K_d}, \qquad
  \Delta\Delta G([\text{NaCl}])
  = \Delta\Delta G_{150}
    \frac{f([\text{NaCl}])}{f(150\,\text{mM})},
  \label{eq:salt_model}
\end{equation}
mit $\Delta\Delta G_{150} = 0{,}8$\,kcal/mol (Tyler et~al.\ 2012,
ZZ-Exchange: $\Delta\Delta G(k_f)=0{,}5$ und $\Delta\Delta G(k_r)=0{,}3$\,kcal/mol,
Summe $0{,}8$).
Damit ist $0{,}8$\,kcal/mol der kodierte 150-mM-Anker und nicht die
Asymptote; die implizierte Asymptote beträgt ungefähr
$0{,}885$\,kcal/mol. Diese Normierung stimmt mit Implementierung und
Tabelle überein. Die genaue Quellensemantik und Unsicherheit benötigen
noch einen Audit des Primärvolltexts.

\begin{table}[!htbp]
\centering
\caption{XCL1-Salzverschiebung: interne Rekonstruktion. Das Modell
ist kalibriert bei $K_\text{eq}(0\,\text{mM}) = 1{,}2$,
$\Delta\Delta G_{150} = 0{,}8$\,kcal/mol und $K_d = 16$\,mM.}
\label{tab:xcl1_salt}
\begin{tabular}{@{}rcccc@{}}
\toprule
{[NaCl]} (mM) & $f(\text{Cl}^-)$ & $K_\text{eq}$ & \% Ltn40 & Experiment \\
\midrule
0   & 0{,}00 & 1{,}200 & 54{,}5 & 54{,}5\,\% (Projektanker) \\
10  & 0{,}39 & 0{,}676 & 40{,}3 & -- \\
50  & 0{,}76 & 0{,}387 & 27{,}9 & -- \\
75  & 0{,}82 & 0{,}351 & 26{,}0 & Tyler: \glqq ähnlich wie $150$\,mM\grqq{}  \\
150 & 0{,}90 & 0{,}311 & 23{,}7 & ${\sim}24\,\%$ Tyler ZZ-exch.\  \\
300 & 0{,}95 & 0{,}291 & 22{,}5 & -- \\
365 & 0{,}96 & 0{,}287 & 22{,}3 & CSP-Sättigung (Tyler) \\
\bottomrule
\end{tabular}
\end{table}

\subsubsection{Unabhängigkeitsstatus}

Die $150$\,mM-Rekonstruktion ($23{,}7\,\%$) verwendet Tylers
$\Delta\Delta G$ als Eingabe und berechnet dieselbe Evidenzfamilie neu.
Sie ist keine unabhängige Vorhersage. Die verbleibenden internen
Kandidatenausgaben sind:

\begin{enumerate}
  \item \textbf{75\,mM-Rasterpunkt:} Das Modell ergibt $26{,}0\,\%$.
    Tylers qualitative Aussage, $70$\,mM und $150$\,mM hätten einen
    ähnlichen Effekt, stammt aus derselben Kalibrierungsevidenz und ist
    kein externer Test.
  \item \textbf{Form der Sättigungskurve:} Die nichtlineare Form wird
    durch Chloridbindungs-Sättigung bestimmt ($K_d = 16$\,mM aus CSP),
    nicht allein aus Tylers Populationsdaten ableitbar.
  \item \textbf{K25A-Mutante:} Die Modellausgabe beträgt
    $K_\text{eq} = 0{,}40$ ($28{,}6\,\%$~Ltn40). Sie wird hier nicht als
    unabhängig vorregistrierte Mutantenvorhersage behandelt.
\end{enumerate}


% --- Result 4 ----------------------------------------------------------------
\subsection[Konstruktionsbedingung: Thermosom]{Konstruktionsbedingung:
Thermosom\\ Meta-Demonstration}
\label{sec:thermosome}

Um zu zeigen, dass PG/non-PG von der Spielformulierung abhängt und nicht
allein durch das biologische System festgelegt ist, analysieren wir das Thermosom von
\emph{Thermoplasma acidophilum} -- ein Gruppe-II-Chaperonin mit
$(\alpha\beta)_4(\alpha\beta)_4$-hexadekamerer Architektur und negativer
Inter-Ring-Kooperativität~\cite{Bigotti2008,Gutsche2000} -- mit zwei bewusst
kontrastierenden Spielkodierungen:

\begin{table}[!htbp]
\centering
\caption{Thermosom: dasselbe biologische System, zwei Spielmodelle,
entgegengesetzte PG-Resultate.}
\label{tab:thermosome}
\footnotesize
\begin{tabularx}{\textwidth}{@{}L{0.22\textwidth}YccrrL{0.13\textwidth}@{}}
\toprule
Modell & Spieler & PG & NE & $I_1$ & $I_2$ & Regime \\
\midrule
A: Ring vs.\ Ring & $\alpha_4\beta_4$-Ring 1 vs.\ Ring 2 &
  $\checkmark$ & 2 & $-12{,}92$ & $-12{,}92$ & GG \\
B: Chaperon vs.\ Substrat & Hexadekamer vs.\ Substrat &
  $\times$ & 1 & $-4{,}0$ & $-1{,}5$ & therm.\ NGG \\
\bottomrule
\end{tabularx}
\end{table}

\textbf{Modell~A} behandelt die zwei Ringe als identische Spieler in
einem symmetrischen Anti-Koordinationsspiel $\{$Tight, Relaxed$\}$. PG
folgt algebraisch, weil $u_1 = u_2^\top$ gilt (identische Spieler
$\Rightarrow I_1 = I_2$ per Konstruktion). Die zwei statischen NE sind
(Tight, Relaxed) und (Relaxed, Tight). Sie stellen komplementäre
Ringzustände dar, belegen aber für sich keine zeitliche Alternation;
dazu wäre ein Übergangs- oder Markov-Modell erforderlich.

\textbf{Modell~B} behandelt den Hexadekamer und das Substrat als
asymmetrische Spieler. Es überführt einen ATP-Bindungseingang und eine
einseitige Substratfaltungs-Beschleunigung in die Scoreterme
$\Delta_H=4{,}0$ und $\Delta_S=1{,}5$. Da diese Observablen für sich keine
gemeinsame thermodynamische Skala definieren, ist ihre Differenz von
$2{,}5$ ein konstruierter PG-Scorekontrast und keine gemessene
Freienergieverletzung.

Das Paar illustriert Modellabhängigkeit und nicht zwei gleichwertig
validierte biologische Mechanismen. Die Wahl ist eine
Modellierungsentscheidung und motiviert die
Konstruktionsbedingungs-Transparenz in der gesamten Arbeit.

\subsubsection{Robustheit}

Ein Scan über $\Delta_H \in [0{,}5,\, 8{,}0] \times \Delta_S \in [0{,}5,\, 8{,}0]$
bestätigt, dass Modell~B PG$=$True ausschließlich auf der Diagonalen
$\Delta_H = \Delta_S$ ergibt, wie algebraisch aus der
Auszahlungsstruktur folgt. Dies ist eine Sensitivität gegenüber dem
festgelegten Scorescan; der Scan belegt nicht, welche
Parameterkombinationen physikalisch plausibel sind.


% =============================================================================
\section{Bekannte Grenzen und Tautologie-Transparenz}
\label{sec:tautology}

\subsection{Was tautologisch ist}

\begin{enumerate}
  \item \textbf{PG/non-PG folgt aus der Auszahlungskonstruktion:}
    $\Delta_H = \Delta_S \Rightarrow$ PG (algebraisch);
    $\Delta_H \neq \Delta_S \Rightarrow$ non-PG (algebraisch).
    Der PG-Test testet die Auszahlungssymmetrie, nicht eine
    Systemeigenschaft.

  \item \textbf{NE-Vielfachheit:} In einem generischen (bindungsfreien)
    exakten $2\times 2$-Potentialspiel gibt es genau dann zwei reine NE,
    wenn kein Spieler eine strikt dominante Strategie besitzt; das ist eine
    Folge der Spielstruktur, kein biologischer Befund.

  \item \textbf{Frühere Nash-Frustration-Resultate (FST-III v3):} Der
    Spektralradius $\rho(I-\eta H)<1$ an einem nichtentarteten Minimum
    mit positiv definiter Hesse-Matrix folgt für
    $0<\eta<2/\lambda_{\max}(H)$. Unter dieser Konstruktion ist er kein
    Diskriminator biologischer Stabilität. Zufallsparameter erfüllen dasselbe Kriterium
    ($\rho(J)<1$ in $50/50$ Zufalls- und $10/10$ gefitteten Läufen).
\end{enumerate}

\noindent\textbf{Skalengrenze.} Exakte Potentialität ist eine kardinale
Eigenschaft. Getrennte positive Skalierungen der Spieler-Utilities
können Beste Antworten und Nash-Gleichgewichte unverändert lassen, aber
$I_H=I_S$ verändern. Vor einer thermodynamischen Interpretation muss
daher eine gemeinsame physikalische Skala hergeleitet werden;
gewichtete und ordinale Potentiale beantworten schwächere, andere
Fragen.

\subsection{Noch nicht unabhängig validierte Kandidatenausgaben}

\begin{enumerate}
  \item \textbf{NE-Identität und Sensitivität:} Die festgelegten
    Matrizen erzeugen Gleichgewichtsprofile und Scoregrenzen, die für
    künftige Tests eingefroren werden können.
  \item \textbf{Salzantwort:} Die kalibrierte XCL1-Rekonstruktion
    erzeugt eine Konzentrationskurve; nach Modell-Freeze wurde jedoch
    noch keine unberührte Konzentrationsmessung geprüft.
  \item \textbf{Regime-Labels:} Die vorgeschlagene Gruppierung ist eine
    nachvollziehbare Rekodierung ausgewählter Symmetriemuster und keine
    empirisch entdeckte Ontologie.
\end{enumerate}

\subsection{Pre-Review-Gate der Validierungsevidenz}

Das maschinenlesbare Validierungsledger aus
\path{scripts/validation_evidence_ledger.py} trennt acht zentrale
Fallrouten entlang von fünf Achsen: Konstruktionsbedingtheit,
experimentelle Dynamik, AlphaFold-Baseline, Baseline der energetischen
Frustration und unabhängiger Chaperon-Zyklusanker. Alle acht Zeilen sind
konstruktionsbedingt; keine enthält eine berechnete AlphaFold-Baseline,
eine Baseline der energetischen Frustration oder einen quantitativen
unabhängigen Zyklusanker. Damit beträgt die Zahl vollständiger
Validierungspässe $0/8$.

Dieses Gate verwirft den engeren XCL1-Populationscheck nicht, erhält aber
seinen Vorbehalt der teilweisen Zirkularität. GroEL bleibt nur ein
qualitativer Teilanker: Zwei von drei Kinetikkategorien sind konsistent,
während sämtliche modellierten Zykluszeiten außerhalb des angegebenen
experimentellen Bereichs von $15$--$20$\,s liegen
(Primärquelle noch zu verifizieren). Der aktuelle Stand
trägt deshalb nur eine diagnostische Klassifikation, kein validiertes
mechanistisches Modell und keine Erweiterung des wissenschaftlichen Anspruchs.


% =============================================================================
\section{Diskussion}
\label{sec:discussion}

\subsection{Ehrliche Bilanz}

Wir bewerten unseren Rahmen an 13~offenen Fragen im Chaperon-Feld
(Tabelle~\ref{tab:balance}). Der Rahmen adressiert ${\sim}38\,\%$ der
Fragen (5/13), davon besitzt nur eine (Q9, XCL1-Salzverschiebung) einen
quantitativen, teilweise zirkulären Konsistenzcheck.

\begin{table}[!htbp]
\centering
\caption{Ehrliche Bilanz: Abdeckung von 13 offenen Fragen durch den
Rahmen.}
\label{tab:balance}
\begin{tabular}{@{}llc@{}}
\toprule
Kategorie & Fragen & Anteil \\
\midrule
QUANTITATIVER CHECK (teilweise zirkulär) & Q9 (XCL1-Salzverschiebung) & 7{,}7\,\% \\
MODELLAUSGABE (konstruktionsbedingt) & Q1, Q6, Q10, Q13 & 30{,}8\,\% \\
RE-DESCRIBE & Q4, Q8 & 15{,}4\,\% \\
NONE & Q2, Q3, Q5, Q7, Q11, Q12 & 46{,}2\,\% \\
\bottomrule
\end{tabular}
\end{table}

${\sim}46\,\%$ der Fragen liegen vollständig außerhalb des Geltungsbereichs
des Rahmens. Diese Bilanz ist das ehrliche Fundament der vorliegenden Arbeit.

\subsection{Vergleich mit bestehenden Ansätzen}

Der PG-Test liefert eine \emph{diagnostische} Klassifikationsschicht, die
bestehende Ansätze ergänzt, aber nicht ersetzt:

\begin{itemize}
  \item \textbf{Statistische Mechanik (Ising, MWC):} Wird das exakte
    Potential eines PG-Spiels mit einer negativen freien Energie auf einer
    gemeinsamen physikalischen Skala identifiziert, reduziert sich die
    Suche nach reinen NE auf das Auffinden von Freienergie-Minima unter
    einseitigen Zügen und fügt der statistischen Mechanik nichts hinzu.
    Innerhalb einer festgelegten kardinalen Repräsentation ist PG dafür
    notwendig, aber nicht hinreichend; die vorliegenden Scores begründen
    keine solche Skala (Abschnitt~\ref{sec:game_construction}).

  \item \textbf{AlphaFold / Strukturvorhersage:} Unser Rahmen arbeitet
    auf einer anderen Ebene -- er kennzeichnet eine festgelegte Spielrepräsentation
    konformationeller Regime, nicht die Struktur selbst. Für metamorphe
    Proteine (XCL1, MAD2) liefert die NE-Analyse
    konstruktionsbedingte Kandidatenbeschreibungen von Bistabilität.

  \item \textbf{Xu-/Goloubinoff-Rahmen:} Xu unterscheidet Mechanismen
    bereits anhand der gebrochenen Symmetriemengen. Unser Beitrag ist
    eine vorgeschlagene spieltheoretische Rekodierung und Prüfung, nicht
    die Entdeckung dieser biologischen Unterscheidung.
\end{itemize}

\subsection{Was die vorliegende Arbeit \emph{nicht} behauptet}

\begin{itemize}
  \item Nash-Gleichgewicht \glqq erklärt\grqq{} Chaperon-Funktion (es
    klassifiziert).
  \item PG/non-PG ist eine \glqq Erklärung\grqq{} für ein Phänomen (es ist
    konstruktionsbedingt).
  \item \glqq Drei Regime statt zwei\grqq{} ist eine Erweiterung von Goloubinoff
    (es ist Unterstruktur).
  \item AF2-Versagen ist durch Nash-Analyse \glqq gelöst\grqq{} (es wird
    umformuliert).
  \item ATP-Kosten werden durch Nash-Analyse \glqq vorhergesagt\grqq{}
    (Goloubinoff erklärt sie bereits).
  \item Die $150$\,mM-Population ist eine unabhängige Vorhersage (sie ist
    eine kalibrierte interne Rekonstruktion).
\end{itemize}


% =============================================================================
\section{Schlussfolgerung}
\label{sec:conclusion}

Der Vierzyklus-Test ist für ein festgelegtes Spiel mathematisch gültig;
seine biologische Interpretation hängt jedoch von Spielerpartition,
Grobkörnung, Nutzendefinition und kardinaler Skala ab. Die
gegenwärtige Evidenz belegt weder eine $\mathrm{S4}\iff\mathrm{PG}$-Äquivalenz
noch einen validierten Atlas mit 16 Systemen oder eine unabhängige
XCL1-Vorhersage. Der vorliegende Beitrag ist ein Konstruktionsledger mit
expliziten Fehlermodi und einem Validierungsstand von $0/8$.

Der nächste entscheidende Schritt für eine eigenständige Evidenz ist eine vorregistrierte
Netzwerk-zu-Spiel-Konstruktion aus gepaarten Vorwärts- und Rückwärtsraten
mit Unsicherheitsfortpflanzung, Hidden-Cycle-Audit und einem verblindeten
externen Endpunkt.

Alle Skripte (36 Python-Dateien) und Resultatdateien sind verfügbar
unter \href{https://github.com/research-line/fst-nash}{GitHub: research-line/fst-nash}.


% =============================================================================
\section*{Daten- und Codeverfügbarkeit}

Alle Berechnungsskripte liegen in \texttt{scripts/} (36~Dateien), die
Resultate in \texttt{results/} (JSON-Format). Die Template-Funktionen für
den Vierzyklus-Test, die Nash-Gleichgewichtssuche und das
Symmetrie-Mapping liegen in \texttt{scripts/fold\_switch\_template.py}.
Hold-out-Skripte liegen in \texttt{scripts/extension\_B/}. Das
Pre-Review-Evidenzgate wird durch
\path{scripts/validation_evidence_ledger.py} erzeugt.


% =============================================================================
\begin{sloppypar}
\section*{KI-Offenlegung}

\noindent
KI-Systeme wurden in erheblichem Ma\ss{}e in allen Phasen dieser Arbeit
eingesetzt: Konzeption und Forschungsdiskussion, Literatur- und
Quellenarbeit, mathematische Analyse und Beweisarbeit, Erstellung und
Pr\"ufung der Rechenskripte, Ausf\"uhrung und Auswertung der
Berechnungsl\"aufe, Text, \"Ubersetzung und Satz.

\medskip

\noindent
Qualit\"atssicherung: Einsatz mehrerer unabh\"angiger Modelle
verschiedener Anbieter mit Gegenpr\"ufungen und adversarialen KI-Reviews;
nachvollziehbare, versionierte Rechenskripte mit gespeicherten
Ergebnisdateien und Pr\"ufsummen; fortlaufende Beweis- und
Registerdokumente; strikte Claim-Hygiene (Behauptungen nur mit
dokumentiertem Beleg, Nachtr\"age statt stiller Korrekturen); Changelogs
\"uber alle Versionen.

\end{sloppypar}


% =============================================================================
\begin{thebibliography}{99}

\bibitem{Xu2022}
H.~Xu,
``Non-equilibrium protein folding and activation by ATP-driven
chaperones,''
\textit{Biomolecules}\ \textbf{12}, 832 (2022), DOI: 10.3390/biom12060832.

\bibitem{Goloubinoff2018}
P.~Goloubinoff, A.~S.~Sassi, B.~Fauvet, A.~Barducci, and P.~De~Los~Rios,
``Chaperones convert the energy from ATP into the nonequilibrium
stabilization of native proteins,''
\textit{Nat.\ Chem.\ Biol.}\ \textbf{14}, 388--395 (2018), DOI: 10.1038/s41589-018-0013-8.

\bibitem{MondererShapley1996}
D.~Monderer and L.~S.~Shapley,
``Potential Games,''
\textit{Games Econ.\ Behav.}\ \textbf{14}, 124--143 (1996), DOI: 10.1006/game.1996.0044.

\bibitem{GeigerFSTIII2025}
L.~Geiger,
``Functional Stability Theory~III: Biological Stability and Nash
Frustration,''
Zenodo preprint (2026), DOI: 10.5281/zenodo.20130573.

\bibitem{Tuinstra2008}
R.~L.~Tuinstra, F.~C.~Peterson, S.~Kutlesa, E.~S.~Elgin,
M.~A.~Kron, and B.~F.~Volkman,
``Interconversion between two unrelated protein folds in the
lymphotactin native state,''
\textit{Proc.\ Natl.\ Acad.\ Sci.\ USA}\ \textbf{105}, 5057--5062 (2008), DOI: 10.1073/pnas.0709518105.

\bibitem{Tyler2012}
R.~C.~Tyler, J.~C.~Wieting, F.~C.~Peterson, and B.~F.~Volkman,
``Electrostatic optimization of the conformational energy landscape in
a metamorphic protein,''
\textit{Biochemistry}\ \textbf{51}, 9067--9075 (2012), DOI: 10.1021/bi300842j.

\bibitem{Bigotti2008}
M.~G.~Bigotti and A.~R.~Clarke,
``Chaperonins: The hunt for the Group~II mechanism,''
\textit{Arch.\ Biochem.\ Biophys.}\ \textbf{474}, 331--339 (2008), DOI: 10.1016/j.abb.2008.03.015.

\bibitem{Gutsche2000}
I.~Gutsche, O.~Mihalache, and W.~Baumeister,
``ATPase cycle of an archaeal chaperonin,''
\textit{J.\ Mol.\ Biol.}\ \textbf{300}, 187--196 (2000), DOI: 10.1006/jmbi.2000.3833.

\end{thebibliography}

\end{document}
